% ============================================================
%  第八章  不完美实验：效应与反事实的定界  （8.4—8.6 及章末）
% ============================================================
\providecommand{\ACE}{\mathrm{ACE}}
\providecommand{\ETT}{\mathrm{ETT}}
\providecommand{\ghist}[6]{%
\begin{scope}[shift={(#3,#4)},scale=#6]
  \draw[->,thin] (-1.15,0)--(1.15,0);
  \foreach \i in {0,...,24}{%
    \pgfmathsetmacro{\xx}{-1+(\i+0.5)*0.08}%
    \pgfmathsetmacro{\hh}{0.95*exp(-((\xx-(#1))*(\xx-(#1)))/(2*(#2)*(#2)))}%
    \fill[black!72] (\xx-0.032,0) rectangle (\xx+0.032,\hh);
  }
\end{scope}
\node[font=\scriptsize] at (#3,#4-0.8) {(#5)};
}

\section{对工具的检验}

如 \S 8.2 所定义，我们的不完美实验模型依赖两个假设：$Z$ 是随机化的，且 $Z$ 对 $Y$ 没有副作用。这两个假设意味着 $Z$ 与 $U$ 独立——经济学家称这一条件为“外生性”（exogeneity），它使 $Z$ 相对于 $X$ 与 $Y$ 之间的关系有资格成为工具变量（见 \S 5.4.3 与 \S 7.4.5）。长期以来，人们认为对“变量 $Z$ 是否外生、是否为工具变量”作实验验证是不可能的（Imbens and Angrist 1994），因为其定义涉及不可观测的因素（或通常所谓的扰动），例如 $U$ 所代表的那些。\footnote{经济学家发展的检验（Wu 1973）只是比较基于两个或更多工具的估计；当出现不一致时，并不能客观地告诉我们哪一个估计不正确。}外生性这一概念，如同因果关系本身，一直被视为主观建模判断的产物，不受非实验数据的检验。

\textup{(8.14a,b)} 给出的界讲述了另一个故事。尽管外生性难以捉摸，它仍可以被赋予一个经验检验。该检验并不保证能检出一切对外生性的违反，但它可以（在某些情形下）筛掉非常糟糕的候选工具。

坚持 \textup{(8.14b)} 中的每个上界都高于 \textup{(8.14a)} 中相应的下界，我们就得到对观测分布的如下可检验约束：
\begin{equation}
\begin{aligned}
P(y_0,x_0\mid z_0)+P(y_1,x_0\mid z_1)&\le 1,\\
P(y_0,x_1\mid z_0)+P(y_1,x_1\mid z_1)&\le 1,\\
P(y_1,x_0\mid z_0)+P(y_0,x_0\mid z_1)&\le 1,\\
P(y_1,x_1\mid z_0)+P(y_0,x_1\mid z_1)&\le 1.
\end{aligned} \tag{8.21}
\end{equation}
若这些不等式中有任何一个被违反，研究者就可以推断：我们模型所依据的假设至少有一个也被违反。如果分配是仔细随机化的，那么对这些不等式的任何违反，都必须归因于分配过程对受试者响应的某种直接影响（例如一次创伤性经历）。反过来，如果 $Z$ 对 $Y$ 的直接效应可以被消除——比如通过有效使用安慰剂——那么任何观测到的违反都可以安全地归因于 $Z$ 与 $U$ 之间的伪相关：即分配偏倚，因而也就是外生性的丧失。Richardson and Robins (2010) 讨论了这些检验的效力。

\subsection*{工具不等式}

\textup{(8.21)} 中的不等式推广到多值变量时，取如下形式
\begin{equation}
\max_{x}\sum_{y}\Bigl[\max_{z}P(y,x\mid z)\Bigr]\le 1, \tag{8.22}
\end{equation}
这被称为\emph{工具不等式}（instrumental inequality）。证明见 Pearl (1995b,c)。把工具不等式推广到 $Z$ 或 $Y$ 连续的情形，并无特别困难。若 $f(y\mid x,z)$ 是给定 $X$ 与 $Z$ 时 $Y$ 的条件密度函数，则不等式变为
\begin{equation}
\int_{y}\max_{z}\bigl[f(y\mid x,z)\,P(x\mid z)\bigr]\,dy\le 1\qquad\forall x. \tag{8.23}
\end{equation}
然而，向连续 $X$ 的过渡标志着行为上的剧变，这使 Pearl (1995c) 猜想：图 8.1 的结构对观测密度不施加任何约束。这一猜想由 Bonet (2001) 证明。

由 \textup{(8.21)} 可见，当控制工具 $Z$ 能在处理 $X$ 保持不变的情况下使响应变量 $Y$ 发生显著变化时，工具不等式就被违反。虽然此类变化原则上可以由 $U$、$X$、$Y$ 之间的强相关来解释（因为 $X$ 并未把 $Z$ 从 $Y$ 中屏蔽掉），但工具不等式为这些变化的幅度设定了限度。

工具不等式与量子物理中的贝尔不等式（Bell's inequality）相似，这并非偶然（Suppes 1988; Cushing and McMullin 1989）；两个不等式都刻画了一类无法用“假设存在潜共同原因”来解释的观测相关。在某种意义上，工具不等式可以视为贝尔不等式在如下情形下的推广：允许相互关联的观测变量 $X$ 与 $Y$ 之间存在直接的因果联系。

如果我们愿意对受试者的行为作附加假设，工具不等式可以被显著收紧——例如假定没有个体会被鼓励工具所劝阻，或者（用数学语言）对一切 $u$ 有
\[
P(x_1\mid z_1,u)\ge P(x_1\mid z_0,u).
\]
这样的假设相当于人群中不存在违抗者，即不存在一贯选择与其分配相反之处理的受试者。在这一假设下，\textup{(8.21)} 的不等式可以被收紧（Balke and Pearl 1997）为
\begin{equation}
\begin{aligned}
P(y,x_1\mid z_1)&\ge P(y,x_1\mid z_0),\\
P(y,x_0\mid z_0)&\ge P(y,x_0\mid z_1),
\end{aligned}\qquad\text{对一切}\ y\in\{y_0,y_1\}. \tag{8.24}
\end{equation}
对这些不等式的违反，现在意味着要么存在选择偏倚，要么 $Z$ 对 $Y$ 有直接效应，要么存在违抗的受试者。

\section{不依从问题的贝叶斯进路}

\subsection{贝叶斯方法与吉布斯抽样}

本节描述一种从有限样本估计因果效应与反事实概率的一般方法，同时纳入关于人群的先验知识。该方法由 Chickering and Pearl (1997) 提出，\footnote{Imbens and Rubin (1997) 提出了一个类似的方法，只是缺少图形视角。}适用于贝叶斯框架：按此框架，(i) 任何未知的统计参数都可以被赋予先验概率；(ii) 对该参数的估计就是计算其后验分布，以抽样得到的数据为条件。在我们的情形中，所涉参数是概率 $P(r_x,r_y)$（简记 $P(r)$），由它可以推出 $\ACE(X\to Y)$。

如果我们不把 $P(r)$ 视为概率，而视为人群中具有 $R=r$ 所给响应特征的个体比例 $v_r$，那么给这样一个量赋予概率的想法，就契合贝叶斯分析的标准哲学：$v_r$ 是一个潜在可测（尽管未知）的物理量，因此可以容许一个先验概率，用以编码我们对该量的不确定性。

假定实验中有 $m$ 名受试者。我们用 $z^i$、$x^i$、$y^i$ 分别表示受试者 $i$ 的 $Z$、$X$、$Y$ 的观测值。类似地，用 $r^i$ 表示受试者 $i$ 的（未观测的）依从（$r_x$）与响应（$r_y$）组合。用 $\chi^i$ 表示三元组 $\{z^i,x^i,y^i\}$。

给定实验中观测到的数据 $\chi$ 以及关于未知比例 $v_r$ 的先验分布，我们的问题是导出 $\ACE(X\to Y)$ 的后验分布。$v_R$ 与 $\ACE(X\to Y)$ 的后验分布都可以用图 8.4 所示的图形模型导出，该模型显式表示了在变量 $\{\chi,v_R,\ACE(X\to Y)\}$ 上定义的联合（贝叶斯）分布中所成立的独立关系。该模型可以理解为响应变量模型（图 8.3）的 $m$ 个实现，每个三元组 $\chi$ 对应一个，并用代表未知比例 $v_R=(v_{r1},v_{r2},\dots,v_{r16})$ 的节点把它们连接起来。该模型显式表示了这样一个假设：给定比例 $v_R$，一名受试者属于 16 个依从—响应子人群中任何一个的概率，不依赖于实验中其他受试者的依从与响应行为。由 \textup{(8.10)}，$\ACE(X\to Y)$ 是 $v_R$ 的确定性函数，因此一旦这些比例已知，$\ACE(X\to Y)$ 就与定义域中一切其他变量独立。

\begin{figure}[htbp]
\centering
\begin{tikzpicture}[causal,scale=1.15,every node/.style={transform shape}]
  \node[vtx] (vR) at (0,4.4) {$v_R$};
  \node[draw,ellipse,inner sep=2.5pt,font=\small] (ACE) at (3.7,4.7) {$\ACE(X\to Y)$};
  \draw[ed] (vR)--(ACE);
  % three subject plates
  \foreach \g/\dx/\tag in {1/-4.0/a,2/-0.8/b,3/2.9/c}{
    \node[vtxs] (\tag1) at (\dx-0.85,3.15) {$z^{\tag}$};
    \node[vtxs] (\tag2) at (\dx+0.65,3.15) {$r^{\tag}$};
    \node[vtxs] (\tag3) at (\dx-0.1,2.2) {$x^{\tag}$};
    \node[vtxs] (\tag4) at (\dx+0.85,1.3) {$y^{\tag}$};
    \draw[ed] (\tag1)--(\tag3);
    \draw[ed] (\tag2)--(\tag3);
    \draw[ed] (\tag3)--(\tag4);
    \draw[ed] (\tag2) to[bend right=28] (\tag4);
    \draw[ed] (vR) to[bend left=8] (\tag2);
  }
  \node[font=\Large] at (1.0,2.5) {$\cdots$};
\end{tikzpicture}
\caption{用于表示 $P(\{\chi\}\cup\{v_R\}\cup\{\ACE(X\to Y)\})$ 中独立关系的模型。}
\label{fig:8.4}
\end{figure}

于是，原则上，估计 $\ACE(X\to Y)$ 就归结为一个完全确定的贝叶斯网络中计算某变量后验概率这一标准推断任务。（用于此推断计算的图形技巧在 \S 1.2.4 中有简要总结。）在许多情形下，图中体现的独立关系可以被利用，使推断任务变得高效。遗憾的是，由于 $r^i$ 从未被观测，在我们的模型中，即便有这些独立关系，导出 $\ACE(X\to Y)$ 的后验分布仍不可处理。为获得 $\ACE(X\to Y)$ 后验分布的估计，可以使用一种称为吉布斯抽样（Gibbs sampling）的近似技巧（Robert and Casella 1999）。该技巧的图形版本被称为“随机模拟”（stochastic simulation），见 Pearl (1988b, p.~210)；其细节（应用于图 8.4 的图时）在 Chickering and Pearl (1997) 中讨论。这里我们以直方图的形式给出典型结果，以展示该技巧对因果推断问题的一般适用性。

\subsection{样本量与先验分布的影响}

该方法的输入有两项：(1) 观测数据 $\chi$，表示为 $\{z,x,y\}$ 的 8 种可能取值各自被观测到的次数；(2) 关于未知比例 $v_R$ 的狄利克雷先验，用 16 个参数表示。系统输出 $\ACE(X\to Y)$ 的后验分布，以直方图形式表示。

为展示先验分布对输出的影响，我们用两种不同的先验给出全部结果。第一种是 16 维向量 $v_R$ 上的平坦（均匀）分布，常用于表达对定义域的无知。第二种先验是有偏的，用以表示受试者的依从特征与响应特征之间存在强依赖。图 8.5 展示了这两种先验分布（在没有任何数据的情况下）所诱导的 $\ACE(X\to Y)$ 的分布。我们看到，图 8.5(b) 的有偏先验把几乎全部权重都赋予了 $\ACE(X\to Y)$ 的负值。

\begin{figure}[htbp]
\centering
\begin{tikzpicture}[causal,every node/.style={transform shape}]
  \ghist{0.00}{0.26}{0}{0}{a}{1.4}
  \ghist{-0.45}{0.22}{4.6}{0}{b}{1.4}
  \node[font=\scriptsize,below] at (-1.55,-0.12) {$-1$};
  \node[font=\scriptsize,below] at (0,-0.12) {$0$};
  \node[font=\scriptsize,below] at (1.55,-0.12) {$1$};
  \node[font=\scriptsize,below] at (3.05,-0.12) {$-1$};
  \node[font=\scriptsize,below] at (4.6,-0.12) {$0$};
  \node[font=\scriptsize,below] at (6.15,-0.12) {$1$};
\end{tikzpicture}
\caption{(a) 由参数 $v_{CR}$ 上的平坦先验所诱导的 $\ACE(X\to Y)$ 的先验分布。(b) 由这些参数上的有偏先验所诱导的 $\ACE(X\to Y)$ 的分布。（原书为吉布斯抽样的输出直方图，此处按原图轮廓重绘为示意，条形高度非逐格复制。）}
\label{fig:8.5}
\end{figure}

为说明样本量的增大会如何冲淡先验分布的影响，我们把该方法应用于从某个分布 $P(x,y\mid z)$ 中抽出的模拟数据，该分布的 ACE 已知是可识别的。这样的分布见表 8.1。对该分布，\textup{(8.14a,b)} 所得的上界与下界收缩为单一的点：
\[
\ACE(X\to Y)=0.55.
\]
图 8.6 展示了吉布斯抽样器应用于从表 8.1 所示分布中抽出的、大小不一的数据集时的输出，分别使用平坦先验与有偏先验。正如预期，随着案例数增加，后验分布越来越集中在 0.55 附近。一般而言，由于 $\ACE(X\to Y)$ 的有偏先验比均匀先验离 0.55 更远，需要更多案例才能使后验分布收敛到 0.55。

\begin{table}[htbp]
\centering
\caption{使 $\ACE(X\to Y)$ 可识别的分布}
\label{tab:8.1}
\small
\begin{tabular}{cccc}
\toprule
$z$ & $x$ & $y$ & $P(x,y,z)$ \\
\midrule
0 & 0 & 0 & 0.275 \\
0 & 0 & 1 & 0.0 \\
0 & 1 & 0 & 0.225 \\
0 & 1 & 1 & 0.0 \\
1 & 0 & 0 & 0.225 \\
1 & 0 & 1 & 0.0 \\
1 & 1 & 0 & 0.0 \\
1 & 1 & 1 & 0.275 \\
\bottomrule
\end{tabular}
\end{table}

\begin{figure}[htbp]
\centering
\begin{tikzpicture}[causal,every node/.style={transform shape}]
  \ghist{0.55}{0.34}{0}{2.6}{a}{1.0}
  \ghist{0.55}{0.24}{3.4}{2.6}{b}{1.0}
  \ghist{0.55}{0.13}{6.8}{2.6}{c}{1.0}
  \ghist{0.55}{0.055}{10.2}{2.6}{d}{1.0}
  \ghist{0.55}{0.28}{0}{0}{e}{1.0}
  \ghist{0.55}{0.18}{3.4}{0}{f}{1.0}
  \ghist{0.55}{0.09}{6.8}{0}{g}{1.0}
  \ghist{0.55}{0.04}{10.2}{0}{h}{1.0}
  \foreach \dx in {0,3.4,6.8,10.2}{
    \node[font=\scriptsize,below] at (\dx+0.55,2.54) {0.55};
    \node[font=\scriptsize,below] at (\dx+0.55,-0.06) {0.55};
  }
\end{tikzpicture}
\caption{使用两种先验时对可识别处理效应的输出直方图：(a)、(b)、(c)、(d) 展示在平坦先验下、数据集分别由 10、100、1000 与 10000 名受试者组成时 $\ACE(X\to Y)$ 的后验；(e)、(f)、(g)、(h) 展示在有偏先验下、相同各数据集时 $\ACE(X\to Y)$ 的后验。（水平线跨区间 $(-1,+1)$。直方图按原图轮廓重绘为示意。）}
\label{fig:8.6}
\end{figure}

\subsection{不完美依从的临床数据中的因果效应}

本节分析两组在不完美依从条件下获得的临床数据。先考虑 \S 8.2.6 所述的脂质研究诊所冠心病一级预防数据。阈值化之后所得的数据集见表 8.2。采用大样本假定，\textup{(8.14a,b)} 给出界
\[
0.39\ \le\ \ACE(X\to Y)\ \le\ 0.78.
\]
图 8.7 展示了基于这些数据的 $\ACE(X\to Y)$ 的后验密度。相当引人注目的是，即便数据集只有 337 个案例，两个后验分布都高度集中在大样本界 0.39 与 0.78 之内。

作为第二个例子，我们考虑 Sommer 等 (1986) 描述的一项实验，该实验旨在确定补充维生素 A 对儿童死亡率的影响。在该研究中，苏门答腊北部的 450 个村庄被随机分配参与一项维生素 A 补充计划，或作为对照组，为期一年。处理组儿童接受两次大剂量维生素 A（$x_1$），对照组儿童不接受处理（$x_0$）。一年期满后，统计两组的死亡人数 $y_0$。该研究的结果也列于表 8.2。

在大样本假定下，\textup{(8.14a,b)} 的不等式给出界
\[
-0.19\ \le\ \ACE(X\to Y)\ \le\ 0.01.
\]
图 8.8 展示了给定数据下、两种先验各自的 $\ACE(X\to Y)$ 的后验密度。有趣的是，对这项研究，先验分布的选择对后验有显著影响。这表明，如果临床医生对自己的先验没有很大把握，就应当作敏感性分析。在这种情况下，渐近界比贝叶斯估计提供的信息更多，而吉布斯抽样的主要作用，将是指示这些界周围边界线的锐利程度。

\begin{table}[htbp]
\centering
\caption{脂质研究与维生素 A 研究的观测数据}
\label{tab:8.2}
\small
\begin{tabular}{ccccc}
\toprule
\multicolumn{3}{c}{} & \multicolumn{2}{c}{观测数} \\
\cmidrule(lr){4-5}
$z$ & $x$ & $y$ & 脂质研究 & 维生素 A 研究 \\
\midrule
0 & 0 & 0 & 158   & 74 \\
0 & 0 & 1 & 14    & 11{,}514 \\
0 & 1 & 0 & 0     & 0 \\
0 & 1 & 1 & 0     & 0 \\
1 & 0 & 0 & 52    & 34 \\
1 & 0 & 1 & 12    & 2{,}385 \\
1 & 1 & 0 & 23    & 12 \\
1 & 1 & 1 & 78    & 9{,}663 \\
\bottomrule
\end{tabular}
\end{table}

\begin{figure}[htbp]
\centering
\begin{tikzpicture}[causal,every node/.style={transform shape}]
  \ghist{0.60}{0.11}{0}{0}{a}{1.5}
  \ghist{0.58}{0.09}{5.4}{0}{b}{1.5}
  \foreach \dx in {0,5.4}{
    \draw[dashed,thin] (\dx+0.39*1.5,0)--(\dx+0.39*1.5,1.5);
    \draw[dashed,thin] (\dx+0.78*1.5,0)--(\dx+0.78*1.5,1.5);
    \node[font=\tiny,below] at (\dx+0.39*1.5,-0.10) {0.39};
    \node[font=\tiny,below] at (\dx+0.78*1.5,-0.34) {0.78};
  }
\end{tikzpicture}
\caption{脂质数据的输出直方图：(a) 使用平坦先验；(b) 使用有偏先验。竖直虚线标出大样本界 0.39 与 0.78。（直方图按原图轮廓重绘为示意。）}
\label{fig:8.7}
\end{figure}

\begin{figure}[htbp]
\centering
\begin{tikzpicture}[causal,every node/.style={transform shape}]
  \ghist{-0.09}{0.09}{0}{0}{a}{1.5}
  \ghist{-0.09}{0.12}{5.4}{0}{b}{1.5}
  \foreach \dx in {0,5.4}{
    \draw[dashed,thin] (\dx-0.19*1.5,0)--(\dx-0.19*1.5,1.5);
    \draw[dashed,thin] (\dx+0.01*1.5,0)--(\dx+0.01*1.5,1.5);
    \node[font=\tiny,below] at (\dx-0.19*1.5,-0.10) {$-0.19$};
    \node[font=\tiny,below] at (\dx+0.01*1.5,-0.34) {0.01};
  }
\end{tikzpicture}
\caption{维生素 A 数据的输出直方图：(a) 使用平坦先验；(b) 使用有偏先验。竖直虚线标出大样本界 $-0.19$ 与 $0.01$。（直方图按原图轮廓重绘为示意。）}
\label{fig:8.8}
\end{figure}

\subsection{单一事件因果的贝叶斯估计}

除评估因果效应外，刚才描述的贝叶斯方法（只需稍作修改）还能回答关于具有特定特征之个体的各种反事实询问。这类询问在 \S 8.3 中已在大样本假定下作了分析与定界。本节演示对如下询问的贝叶斯分析：已知 (1) Joe 在脂质研究的对照组；(2) Joe 按规定服用了安慰剂；(3) Joe 的胆固醇水平没有改善，那么若 Joe 服用消胆胺，他的胆固醇读数得到改善的概率是多少？

这一询问可以通过在与图 8.4 相同的模型上运行吉布斯抽样器来回答，只需把函数 $\ACE(X\to Y)$（式 \textup{(8.10)}）替换为 $v_R$ 的另一个函数，即表示我们询问的函数。若 Joe 在对照组并服用了安慰剂，则他要么是依从者，要么是永不接受者。此外，由于 Joe 的胆固醇水平没有改善，Joe 的响应行为要么是永不康复，要么是受益。因此，他必属于以下四个依从—响应人群之一：
\[
\{(r_x=0,r_y=0),\ (r_x=0,r_y=1),\ (r_x=1,r_y=0),\ (r_x=1,r_y=1)\}.
\]
若 Joe 服用消胆胺，他本会改善，当且仅当其响应行为是受益（$r_y=1$）。由此，所关心的询问由如下函数刻画：
\[
f(v_R)=\frac{v_{01}+v_{11}}{v_{00}+v_{01}+v_{10}+v_{11}}.
\]
图 8.9(a) 与 (b) 分别展示了平坦先验与有偏先验所导致的 $f(v_R)$ 的先验分布。图 8.9(c) 与 (d) 分别展示了使用平坦先验与有偏先验时、由脂质数据得到的 $f(v_R\mid\chi)$ 的后验分布。作为参照，大样本假定下算得的界为
\[
0.51\ \le\ f(v_R\mid\chi)\ \le\ 0.86.
\]
因此，尽管处理组中有 39\% 的不依从，尽管只有 337 名受试者，该研究仍有力地支持这一结论：给定 Joe 的具体经历，他服用该药本会更好。而且，这一结论对两种先验都成立。

\begin{figure}[htbp]
\centering
\begin{tikzpicture}[causal,every node/.style={transform shape}]
  \ghist{0.62}{0.16}{0}{3.4}{a}{1.3}
  \ghist{0.70}{0.12}{4.6}{3.4}{b}{1.3}
  \ghist{0.70}{0.09}{0}{0}{c}{1.3}
  \ghist{0.72}{0.08}{4.6}{0}{d}{1.3}
  \foreach \dx in {0,4.6}{
    \draw[dashed,thin] (\dx+0.51*1.3,0)--(\dx+0.51*1.3,1.3);
    \draw[dashed,thin] (\dx+0.86*1.3,0)--(\dx+0.86*1.3,1.3);
    \draw[dashed,thin] (\dx+0.51*1.3,3.4)--(\dx+0.51*1.3,4.7);
    \draw[dashed,thin] (\dx+0.86*1.3,3.4)--(\dx+0.86*1.3,4.7);
    \node[font=\tiny,below] at (\dx+0.51*1.3,-0.10) {0.51};
    \node[font=\tiny,below] at (\dx+0.86*1.3,-0.34) {0.86};
  }
\end{tikzpicture}
\caption{由反事实询问“若 Joe 服用了该药、而他在未服药时并未改善，他本会改善吗？”所指定的子人群 $f(v_R)$ 的先验（(a) 与 (b)）与后验（(c) 与 (d)）分布。(a) 对应平坦先验，(b) 对应有偏先验。竖直虚线标出大样本界 0.51 与 0.86。（直方图按原图轮廓重绘为示意。）}
\label{fig:8.9}
\end{figure}

\section{结论}

本章发展了因果分析技巧，以处理临床试验中的一大难题：在不完美依从的情况下评估处理效力。仅基于意向性治疗分析的估计——以及基于工具变量公式的估计——都可能产生误导，甚至可能完全落在理论界之外。本章建立的一般公式为政策分析提供了不依赖于工具的保证；此外，它们应当能使分析者判定：为强化依从而作的努力，可以在多大程度上提高总体处理效力。

间接实验与工具变量的重要性并不限于以人为受试者的研究。只要我们试图评估其因果效应的变量无法被直接操纵、却可以被间接手段部分影响，就会出现与不完美依从相当的实验条件。典型的应用涉及对进行中的过程作诊断：其中故障行为的源头必须借助间接手段来识别，因为对可疑源头的直接操纵或者在物理上不可能，或者代价高得难以承受。

在方法论上，本章传达的信息是：即使在因果量不可识别的情形下，关于该领域中因果关系结构的合理假设，也可以被用来获得关于这些关系强度的有用定量信息。一旦这样的假设以图形形式被表述、并以规范划分重新编码，它们就可以交给代数方法处理，从而得到所关心之量的有信息量的界。规范划分还使我们能够用关于所研究人群的先验信念来补充结构假设，并借助吉布斯抽样技巧来便利目标量的贝叶斯估计。

\section*{致谢}
\addcontentsline{toc}{section}{致谢}

本章基于与 Alex Balke 和 David Chickering 的合作工作，并得益于与 Phil Dawid、James Heckman、Guido Imbens、Steve Jacobsen、Steffen Lauritzen、Charles Manski 与 James Robins 的讨论。
